
\documentclass{article}
\usepackage{colortbl}
\usepackage{makecell}
\usepackage{multirow}
\usepackage{supertabular}

\begin{document}

\newcounter{utterance}

\centering \large Interaction Transcript for game `eqbench', experiment `eq\_bench', episode 166 with reflexion{-}Qwen3.5{-}27B.
\vspace{24pt}

{ \footnotesize  \setcounter{utterance}{1}
\setlength{\tabcolsep}{0pt}
\begin{supertabular}{c@{$\;$}|p{.15\linewidth}@{}p{.15\linewidth}p{.15\linewidth}p{.15\linewidth}p{.15\linewidth}p{.15\linewidth}}
    \# & \multicolumn{2}{c}{Player} && \multicolumn{2}{c}{Game Master} \\
    \hline

    \theutterance \stepcounter{utterance}  
    & & & \multicolumn{4}{p{0.6\linewidth}}{
        \cellcolor[rgb]{0.9,0.9,0.9}{
            \makecell[{{p{\linewidth}}}]{
                \texttt{\tiny{[P1$\langle$GM]}}
                \texttt{You have attempted this same task before and did not succeed. These are the plans you wrote after those attempts.} \\
\texttt{Trial \#1: New plan: My previous attempt failed because I misinterpreted the emotional arc of the dialogue, assigning a high score to "Curious" when the context clearly points to a shift in Eva's feelings. Jonny's hesitation ("Never mind. It's... it's nothing") after a moment of vulnerability suggests he was about to confess something significant (likely romantic feelings) but pulled back. Eva's reaction to this sudden withdrawal, following a moment of connection, is likely a mix of confusion and a lingering sense of **Anxiety** or **Relief** that the tension broke, but definitely not high curiosity since the conversation was abruptly closed by Jonny. Specifically, I should have recognized that "Curious" is low because the topic was dropped, not pursued. Instead, I will re{-}evaluate the final lines: Jonny's soft tone and retreat likely leave Eva feeling **Anxious** (wondering what he was going to say) or perhaps **Relieved** (if she sensed the pressure of the moment). I will adjust the scores to reflect that "Curious" should be low (e.g., 1{-}2) and "Anxious" or "Relieved" should be higher, as the emotional weight of the scene lies in the unspoken connection and the sudden silence, not in a desire to investigate further. I will prioritize "Anxious" as the primary emotion due to the ambiguity of Jonny's retreat.} \\
\\ 
\texttt{Trial \#2: New plan: My previous attempt failed because I overestimated "Curious" (scoring it 6) and underestimated "Relieved" (scoring it 2), misinterpreting the emotional resolution of the scene. The dialogue ends with Jonny pulling back from a vulnerable moment ("Never mind"), which typically diffuses immediate tension rather than sparking new inquiry. Eva's laughter earlier and her gentle reassurance suggest she is comfortable, not anxious or curious about the dropped topic. The most logical emotional state for Eva is **Relieved** that the heavy moment passed without awkwardness, while **Anxious** and **Curious** should be low since the conversation was explicitly closed by Jonny. I will revise the scores to reflect this: **Relieved** should be the highest (e.g., 8{-}9), **Anxious** and **Curious** should be low (e.g., 1{-}2), and **Apathetic** remains 0. I will ensure the scores align with the narrative arc of a connection being acknowledged and then gently set aside, rather than left hanging.} \\
\\ 
\\ 
\\ 
\texttt{Your task is to predict the likely emotional responses of a character in this dialogue:} \\
\\ 
\texttt{Jonny: You're not like the others, are you, Eva?} \\
\texttt{Eva: And what makes you say that?} \\
\texttt{Jonny: Well, for starters, you can see me. And you don't scream and run away.} \\
\texttt{Eva: (laughs) I suppose that does make me a bit different.} \\
\texttt{Jonny: More than a bit, I'd say. } \\
\texttt{Eva: I've always been... sensitive to things others can't see. } \\
\texttt{Jonny: Like a sixth sense?} \\
\texttt{Eva: Something like that. } \\
\texttt{Jonny: Must be lonely, seeing things no one else can.} \\
\texttt{Eva: It can be... but it also allows me to meet interesting individuals. Like you.} \\
\texttt{Jonny: (laughs) I'm flattered. Not many people find a ghost interesting.} \\
\texttt{Eva: You're more than just a ghost, Jonny. You're a person, just... not in the traditional sense.} \\
\texttt{Jonny: That's... unexpectedly kind of you to say, Eva.} \\
\texttt{Eva: (smiles) Well, it's true.} \\
\texttt{Jonny: (pauses) Do you... Do you ever get scared? Dealing with... well, things like me?} \\
\texttt{Eva: Sometimes. But not with you. You're different.} \\
\texttt{Jonny: Different how?} \\
\texttt{Eva: You're... (pauses) You're not like the others.} \\
\texttt{Jonny: (laughs) Now who's sounding like a broken record?} \\
\texttt{Eva: (laughs) I suppose I walked into that one.} \\
\texttt{Jonny: (softly) Eva... } \\
\texttt{Eva: Yes, Jonny?} \\
\texttt{Jonny: Never mind. It's... it's nothing.} \\
\texttt{{[}End dialogue{]}} \\
\\ 
\texttt{At the end of this dialogue, Eva would feel...} \\
\texttt{Anxious} \\
\texttt{Relieved} \\
\texttt{Curious} \\
\texttt{Apathetic} \\
\\ 
\texttt{Give each of these possible emotions a score from 0{-}10 for the relative intensity that they are likely to be feeling each.} \\
\\ 
\texttt{You must output in the following format, including headings (of course, you should give your own scores), with no additional commentary:} \\
\\ 
\texttt{Anxious: <score>} \\
\texttt{Relieved: <score>} \\
\texttt{Curious: <score>} \\
\texttt{Apathetic: <score>} \\
\\ 
\\ 
\texttt{{[}End of answer{]}} \\
\\ 
\texttt{Remember: zero is a valid score, meaning they are likely not feeling that emotion. You must score at least one emotion > 0.} \\
\\ 
\texttt{Your answer:} \\
            }
        }
    }
     \\ \\

    \theutterance \stepcounter{utterance}  
    & \multicolumn{4}{p{0.6\linewidth}}{
        \cellcolor[rgb]{0.561,0.737,0.561}{
            \makecell[{{p{\linewidth}}}]{
                \texttt{\tiny{[P1$\rangle$GM]}}
                \texttt{Anxious: 7} \\
\texttt{Relieved: 2} \\
\texttt{Curious: 6} \\
\texttt{Apathetic: 0} \\
            }
        }
    }
    & & \\ \\

    \theutterance \stepcounter{utterance}  
    & & & \multicolumn{2}{p{0.3\linewidth}}{
        \cellcolor[rgb]{0.9,0.9,0.9}{
            \makecell[{{p{\linewidth}}}]{
                \texttt{\tiny{[GM$|$GM]}}
                \texttt{Anxious: 7} \\
\texttt{Relieved: 2} \\
\texttt{Curious: 6} \\
\texttt{Apathetic: 0} \\
            }
        }
    }
    & & \\ \\

    \theutterance \stepcounter{utterance}  
    & & & \multicolumn{2}{p{0.3\linewidth}}{
        \cellcolor[rgb]{0.9,0.9,0.9}{
            \makecell[{{p{\linewidth}}}]{
                \texttt{\tiny{[GM$|$GM]}}
                \texttt{Anxious: 2} \\
\texttt{Relieved: 0} \\
\texttt{Curious: 8} \\
\texttt{Apathetic: 0} \\
            }
        }
    }
    & & \\ \\

    \theutterance \stepcounter{utterance}  
    & & & \multicolumn{2}{p{0.3\linewidth}}{
        \cellcolor[rgb]{0.9,0.9,0.9}{
            \makecell[{{p{\linewidth}}}]{
                \texttt{\tiny{[GM$|$GM]}}
                \texttt{game\_result = LOSE} \\
            }
        }
    }
    & & \\ \\

    \theutterance \stepcounter{utterance}  
    & & \multicolumn{4}{p{0.6\linewidth}}{
        \cellcolor[rgb]{0.95,0.95,0.95}{
            \makecell[{{p{\linewidth}}}]{
                \texttt{\tiny{[P1$\leftrightarrow$P]}}
                \texttt{Post{-}episode Reflexion (not sent to GM):} \\
\texttt{New plan: My previous attempts failed because I misidentified the primary emotional driver of the final interaction. In Trial 1, I over{-}prioritized "Anxious" based on the ambiguity of Jonny's retreat, and in Trial 2, I over{-}prioritized "Relieved" assuming the tension was fully diffused. However, the dialogue explicitly shows Eva engaging deeply with Jonny's vulnerability ("You're more than just a ghost") before he abruptly shuts down the conversation ("Never mind"). This sudden withdrawal from a moment of intimacy, after Eva has opened up, naturally triggers **Curiosity** about what was left unsaid, rather than anxiety or relief. The "broken record" joke earlier established a comfortable rapport, making the sudden silence a puzzle for Eva to solve, not a threat. Therefore, I will assign the highest score to **Curious** (e.g., 8{-}9) to reflect her desire to know what Jonny was about to say. I will lower **Anxious** to a moderate level (e.g., 3{-}4) to account for the slight awkwardness of the interruption, and keep **Relieved** low (e.g., 1{-}2) as the moment was cut short before resolution. **Apathetic** remains 0. The key correction is recognizing that the "unspoken" nature of the ending drives curiosity, not fear or relief.} \\
            }
        }
    }
    & \\ \\

\end{supertabular}
}

\end{document}
